\documentclass[11pt,a4paper]{article}

\usepackage[margin=1in]{geometry}
\usepackage{amsmath,amssymb,amsthm}
\usepackage{booktabs}
\usepackage{array}
\usepackage[hidelinks]{hyperref}

\theoremstyle{plain}
\newtheorem{theorem}{Theorem}[section]
\newtheorem{lemma}[theorem]{Lemma}
\newtheorem{proposition}[theorem]{Proposition}
\newtheorem{corollary}[theorem]{Corollary}

\theoremstyle{definition}
\newtheorem{definition}[theorem]{Definition}
\newtheorem{remark}[theorem]{Remark}

\newcommand{\CDinf}{\mathrm{CD}_{\wedge}}
\newcommand{\CDsup}{\mathrm{CD}_{\vee}}
\newcommand{\op}{\mathrm{op}}
\newcommand{\DM}{\mathrm{DM}}
\newcommand{\Down}{\mathcal{O}}

\title{\textbf{Disproof and Complete Structural Analysis of Conjecture 00000008559:\\
Complete Distributivity, Chain Products, and Minimal Counterexamples}}
\author{Formal Mathematical Report \& Lean 4 Verification}
\date{November 2026}

\begin{document}
\maketitle

\begin{abstract}
We resolve Conjecture \texttt{00000008559} in the negative (\textbf{disproved}).
The conjecture asserts the conjunction of two claims about complete lattices:
(1) a complete lattice is completely distributive if and only if it is ``generated by
complete chain products,'' and (2) ``the smallest complete lattice failing complete
distributivity is the dual of the three-element antichain construction.''
We establish an unconditional disproof of the second clause under every standard
order-theoretic reading of both ``smallest'' (minimum cardinality with uniqueness up to
order isomorphism, or order/sublattice embeddability into every counterexample) and
``the dual of the three-element antichain construction'' (bounded/Dedekind--MacNeille
completion $\DM(A_3)^{\op} = M_3^{\op} \cong M_3$, or down-set completion
$\Down(A_3)^{\op} \cong B_3^{\op}$). Specifically, we prove that every complete lattice
of cardinality $|L| \le 4$ is completely distributive, while at the sharp minimum
cardinality $|L| = 5$ there exist \emph{two} non-isomorphic complete lattices failing
complete distributivity---the modular diamond lattice $M_3$ and the non-modular pentagon
lattice $N_5$---neither of which order-embeds into the other; meanwhile
$\Down(A_3)^{\op} \cong B_3^{\op}$ has cardinality $8$ and \emph{satisfies} complete
distributivity. For the first clause, we prove Raney's positive subdirect theorem that
every complete sublattice of a direct product of complete chains is completely
distributive, while also proving that literal internal complete-lattice generation and
complete-join quotient generation by complete chain products both fail on $M_3$.
All results are formalized from first principles in Core Lean 4 (\texttt{v4.19.0})
with zero external dependencies, zero \texttt{sorry}, and zero \texttt{native\_decide}.
\end{abstract}

\section{Official Problem Statement and Verdict}

Conjecture \texttt{00000008559} reads verbatim as follows:

\begin{quote}
\textbf{English.} Definition: The infinite distributive law for complete lattices.
Conjecture: A complete lattice is completely distributive if and only if it is generated
by complete chain products; and the smallest complete lattice failing complete
distributivity is the dual of the three-element antichain construction.
(\emph{complete distributivity chain product characterization})
\end{quote}

The accompanying Chinese text translates literally as: \emph{``Definition: The infinite
distributive law of complete lattices. Conjecture: A complete lattice satisfies complete
distributivity iff it can be generated by complete chain products; the smallest complete
lattice that is not completely distributive is the dual of the three-antichain.''}

\medskip
\noindent\textbf{Verdict: DISPROVED (FALSE).}
Conjecture \texttt{00000008559} is a logical conjunction $P_1 \wedge P_2$:
\begin{itemize}
  \item \textbf{Clause 2 ($P_2$) is unconditionally false under all standard readings.}
    Every complete lattice with $|L| \le 4$ is completely distributive
    (Theorem~\ref{thm:size4}). At the minimal counterexample cardinality $|L| = 5$,
    there is \emph{no unique} complete lattice failing complete distributivity: both the
    $5$-element diamond lattice $M_3 = \hat{0} \oplus A_3 \oplus \hat{1}$ (which is
    moreover self-dual, $M_3^{\op} \cong M_3$) and the $5$-element pentagon lattice $N_5$
    are complete lattices failing complete distributivity, and $M_3^{\op} \cong M_3 \not\cong N_5$
    because $M_3$ is modular whereas $N_5$ is non-modular (Proposition~\ref{prop:M3-N5}).
    Neither order-embeds into the other, so no single complete lattice $L_*$ is smallest
    either by cardinality (up to isomorphism) or by order/sublattice embedding
    (Theorem~\ref{thm:clause2-false}). Under the alternative reading of ``the dual of the
    three-element antichain construction'' as the dual of the down-set (Alexandroff /
    order-ideal) completion $\Down(A_3)^{\op} \cong (C_2^3)^{\op} = B_3^{\op}$, the
    resulting lattice has $|B_3^{\op}| = 8 > 5$ and \emph{is} completely distributive
    (Proposition~\ref{prop:downset-A3}).
  \item \textbf{Clause 1 ($P_1$) is false under literal algebraic/order-theoretic
    generation, and holds only when repaired to Raney's complete-sublattice
    characterization.} By Raney's representation theorem~\cite{Raney1953}, a complete
    lattice $L$ is completely distributive if and only if $L$ embeds as a \emph{complete
    sublattice} into a direct product of complete chains (we prove the sufficient
    direction in full generality in Theorem~\ref{thm:raney-pos}). However, under the
    literal wording ``generated by complete chain products,'' $M_3$ is both the union of
    three complete sublattices isomorphic to the complete chain $C_3$ (internal
    complete-lattice generation) and a surjective complete-join homomorphic image of the
    complete chain product $C_2 \times C_2 \times C_2$ (closure-quotient generation), yet
    $M_3$ fails complete distributivity (Theorem~\ref{thm:clause1-literal}).
  \item Because $P_2$ is false for \emph{every} candidate complete lattice $L_*$, the
    conjunction $P_1 \wedge P_2$ is false regardless of whether $P_1$ is read literally
    or charitably repaired via Raney's theorem (Theorem~\ref{thm:master-disproof}).
\end{itemize}

\section{Complete Lattices and the Infinite Distributive Law}

\begin{definition}[Complete lattice]
\label{def:complete-lattice}
A partially ordered set $(L, \le)$ is a \emph{complete lattice} if every subset
$S \subseteq L$ admits both a greatest lower bound (infimum) $\bigwedge S \in L$ and a
least upper bound (supremum) $\bigvee S \in L$, characterized for all $z \in L$ by
\[
  z \le \bigwedge S \iff (\forall x \in S,\; z \le x),
  \qquad
  \bigvee S \le z \iff (\forall x \in S,\; x \le z).
\]
In particular, $\bot := \bigvee \emptyset$ and $\top := \bigwedge \emptyset$ satisfy
$\bot \le x \le \top$ for all $x \in L$, and binary meet and join are given by
$x \wedge y := \bigwedge \{x, y\}$ and $x \vee y := \bigvee \{x, y\}$. For any indexed
family $\{y_i\}_{i \in I}$ in $L$, we write $\bigwedge_{i \in I} y_i := \bigwedge \{y_i : i \in I\}$
and $\bigvee_{i \in I} y_i := \bigvee \{y_i : i \in I\}$.
\end{definition}

For any complete lattice $L$, its \emph{order dual} $L^{\op}$ has the same underlying set
with reversed order $x \le_{\op} y \iff y \le x$, swapping $\bigwedge^{\op} = \bigvee$
and $\bigvee^{\op} = \bigwedge$. For any family $\{L_\alpha\}_{\alpha \in A}$ of complete
lattices, the direct product $\prod_{\alpha \in A} L_\alpha$ is a complete lattice under
coordinatewise order, infima, and suprema. An \emph{order isomorphism} $e : L \xrightarrow{\sim} M$
is a bijection satisfying $x \le y \iff e(x) \le e(y)$ for all $x, y \in L$; every order
isomorphism automatically preserves arbitrary infima and suprema.

\begin{definition}[Infinite distributive laws and complete distributivity]
\label{def:cd}
Let $L$ be a complete lattice.
\begin{enumerate}
  \item $L$ satisfies the \emph{meet-over-join infinite distributive law} ($\CDinf$) if
    for every index set $I$, every family of index sets $\{J(i)\}_{i \in I}$, and every
    doubly-indexed family $\{x_{i,j}\}_{i \in I,\, j \in J(i)}$ in $L$,
    \begin{equation}
      \label{eq:cd-inf-sup}
      \bigwedge_{i \in I} \bigvee_{j \in J(i)} x_{i,j}
      \;=\;
      \bigvee_{f \in \prod_{i \in I} J(i)} \bigwedge_{i \in I} x_{i, f(i)}.
    \end{equation}
  \item $L$ satisfies the \emph{join-over-meet infinite distributive law} ($\CDsup$) if
    for every $\{x_{i,j}\}_{i \in I,\, j \in J(i)}$ in $L$,
    \begin{equation}
      \label{eq:cd-sup-inf}
      \bigvee_{i \in I} \bigwedge_{j \in J(i)} x_{i,j}
      \;=\;
      \bigwedge_{f \in \prod_{i \in I} J(i)} \bigvee_{i \in I} x_{i, f(i)}.
    \end{equation}
  \item $L$ is \emph{completely distributive} if both $(\CDinf)$ and $(\CDsup)$ hold.
\end{enumerate}
\end{definition}

Note that no nonemptiness restrictions on $I$ or $J(i)$ are required: when $I = \emptyset$,
both sides of~\eqref{eq:cd-inf-sup} equal $\top$ (since $\prod_{i \in \emptyset} J(i)$ is a
singleton); when $J(i_0) = \emptyset$ for some $i_0 \in I$, both sides equal $\bot$
(since $\prod_{i \in I} J(i) = \emptyset$). Moreover, in~\eqref{eq:cd-inf-sup} the
inequality $\mathrm{RHS} \le \mathrm{LHS}$ holds in \emph{every} complete lattice, since
$\bigwedge_{i'} x_{i', f(i')} \le x_{i, f(i)} \le \bigvee_{j} x_{i,j}$ for all $f$ and $i$;
dually, $\mathrm{LHS} \le \mathrm{RHS}$ in~\eqref{eq:cd-sup-inf} holds universally.

\begin{proposition}[Raney's Self-Duality Theorem~\cite{Raney1952}]
\label{prop:self-duality}
In every complete lattice $L$, $(\CDinf)$ holds if and only if $(\CDsup)$ holds.
Consequently, $L$ is completely distributive if and only if $L^{\op}$ is completely
distributive.
\end{proposition}

\begin{proof}
Suppose $(\CDinf)$ holds in $L$, and let $\{x_{i,j}\}_{i \in I,\, j \in J(i)}$ be given
with $F := \prod_{i \in I} J(i)$. It suffices to prove
$\bigwedge_{f \in F} \bigvee_{i \in I} x_{i, f(i)} \le \bigvee_{i \in I} \bigwedge_{j \in J(i)} x_{i,j}$.
Applying $(\CDinf)$ to the family $y_{f,i} := x_{i, f(i)}$ indexed by $f \in F$ and $i \in I$ gives
\[
  \bigwedge_{f \in F} \bigvee_{i \in I} x_{i, f(i)}
  \;=\;
  \bigvee_{\Phi \in I^F} \bigwedge_{f \in F} x_{\Phi(f),\, f(\Phi(f))}.
\]
Fix $\Phi : F \to I$, and set $u_\Phi := \bigwedge_{f \in F} x_{\Phi(f),\, f(\Phi(f))}$.
We claim $u_\Phi \le \bigvee_{i \in I} \bigwedge_{j \in J(i)} x_{i,j}$. Suppose otherwise;
then for every $i \in I$, we have $u_\Phi \not\le \bigwedge_{j \in J(i)} x_{i,j}$, so
by the Axiom of Choice there exists $f^* \in F$ with $u_\Phi \not\le x_{i, f^*(i)}$ for
every $i \in I$. Taking $i = \Phi(f^*)$ yields $u_\Phi \not\le x_{\Phi(f^*),\, f^*(\Phi(f^*))}$,
contradicting the definition of $u_\Phi$ as the infimum over all $f \in F$. Taking the
supremum over $\Phi \in I^F$ establishes $(\CDsup)$. For the converse, if $(\CDsup)$ holds
in $L$, then $(\CDinf)$ holds in $L^{\op}$, hence $(\CDsup)$ holds in $L^{\op}$, which is
$(\CDinf)$ in $L$.
\end{proof}

\begin{proposition}[Subset collapse and universe invariance]
\label{prop:subset-collapse}
If $(\CDinf)$ holds for index families bounded in cardinality by $|\mathcal{P}(L)|$
(in type theory, index types in the same universe $\mathrm{Type}\,u$ as $L$), then
$(\CDinf)$ holds for arbitrary index families $\{x_{i,j}\}_{i \in I,\, j \in J(i)}$ in any
universes, and in particular implies ordinary finitary distributivity
$x \wedge (y \vee z) = (x \wedge y) \vee (x \wedge z)$.
\end{proposition}

\begin{proof}
For each $i \in I$, put $S_i := \{x_{i,j} : j \in J(i)\} \subseteq L$, and let
$\mathcal{A} := \{S_i : i \in I\} \subseteq \mathcal{P}(L)$. Then
$\bigwedge_{i \in I} \bigvee_{j \in J(i)} x_{i,j} = \bigwedge_{S \in \mathcal{A}} \bigvee_{s \in S} s
= \bigvee_{g \in \prod_{S \in \mathcal{A}} S} \bigwedge_{S \in \mathcal{A}} g(S)$.
For each choice $g \in \prod_{S \in \mathcal{A}} S$ and each $i \in I$, since $g(S_i) \in S_i$,
we may choose $f(i) \in J(i)$ with $x_{i, f(i)} = g(S_i)$, giving
$\bigwedge_{S \in \mathcal{A}} g(S) \le \bigwedge_{i \in I} x_{i, f(i)}$. Taking suprema
proves~\eqref{eq:cd-inf-sup}. Specializing to $I = \{0, 1\}$, $J(0) = \{0\}$ with
$x_{0,0} = x$, and $J(1) = \{0, 1\}$ with $x_{1,0} = y, x_{1,1} = z$ yields
$x \wedge (y \vee z) = (x \wedge y) \vee (x \wedge z)$.
\end{proof}

\begin{theorem}[Complete distributivity on finite complete lattices]
\label{thm:finite-cd}
Let $L$ be a finite complete lattice. Then $L$ is completely distributive (over all,
possibly infinite, index families) if and only if $L$ is distributive, i.e.,
$x \wedge (y \vee z) = (x \wedge y) \vee (x \wedge z)$ for all $x, y, z \in L$.
\end{theorem}

\begin{proof}
Necessity follows from Proposition~\ref{prop:subset-collapse}. Conversely, suppose $L$
is finite and satisfies $x \wedge (y \vee z) = (x \wedge y) \vee (x \wedge z)$. Since
$x \wedge \bot = \bot$, induction on finite lists shows first that
$a \wedge \bigvee_{b \in B} b = \bigvee_{b \in B} (a \wedge b)$ for any finite list $B$,
then $\bigvee_{a \in A} a \wedge \bigvee_{b \in B} b = \bigvee_{(a,b) \in A \times B} (a \wedge b)$,
and by induction on a finite list of finite lists $S_1, \dots, S_m$ in $L$ that
\[
  \bigwedge_{k=1}^m \bigvee S_k
  \;=\;
  \bigvee_{(s_1, \dots, s_m) \in S_1 \times \cdots \times S_m} (s_1 \wedge \cdots \wedge s_m).
\]
Now let $\{x_{i,j}\}_{i \in I,\, j \in J(i)}$ be \emph{any} (possibly infinite) doubly-indexed
family in $L$. Because $L$ is finite, each row set $S_i := \{x_{i,j} : j \in J(i)\} \subseteq L$
is finite, and the collection $\{S_i : i \in I\} \subseteq \mathcal{P}(L)$ is a finite set
$\{R_1, \dots, R_m\}$ of at most $2^{|L|}$ distinct subsets of $L$. Thus
$\bigwedge_{i \in I} \bigvee_{j \in J(i)} x_{i,j} = \bigwedge_{k=1}^m \bigvee R_k$.
Applying the finite-product identity to $R_1, \dots, R_m$, every tuple
$c = (r_1, \dots, r_m) \in R_1 \times \cdots \times R_m$ picks for each $i \in I$ an element
$r_{k(i)} \in R_{k(i)} = S_i$, hence some $f(i) \in J(i)$ with $\bigwedge c \le r_{k(i)} = x_{i, f(i)}$.
Therefore $\bigwedge c \le \bigwedge_{i \in I} x_{i, f(i)}$, establishing $(\CDinf)$, and
$(\CDsup)$ follows by Proposition~\ref{prop:self-duality}.
\end{proof}

\section{Clause 1: Chain Products vs.\ Literal Generation}

A complete lattice $C$ is a \emph{complete chain} if its order is total ($x \le y$ or
$y \le x$ for all $x, y \in C$). A map $\varphi : L \to P$ between complete lattices is a
\emph{complete-lattice embedding} (making $L$ isomorphic to a \emph{complete sublattice}
of $P$) if $\varphi$ is injective and preserves arbitrary infima and suprema:
$\varphi(\bigwedge S) = \bigwedge \varphi[S]$ and $\varphi(\bigvee S) = \bigvee \varphi[S]$
for all $S \subseteq L$.

\begin{theorem}[Positive chain-product theorem; cf.\ Raney~\cite{Raney1953}]
\label{thm:raney-pos}
\begin{enumerate}
  \item Every complete chain $C$ is completely distributive.
  \item Every direct product $\prod_{\alpha \in A} L_\alpha$ of completely distributive
    complete lattices is completely distributive.
  \item Every complete lattice $L$ admitting a complete-lattice embedding into a direct
    product of complete chains $\prod_{\alpha \in A} C_\alpha$ is completely distributive.
\end{enumerate}
\end{theorem}

\begin{proof}
(1) Let $\{x_{i,j}\}_{i \in I,\, j \in J(i)}$ be a family in a complete chain $C$, and put
$\ell := \bigwedge_{i \in I} \bigvee_{j \in J(i)} x_{i,j}$ and
$r := \bigvee_{f \in \prod_i J(i)} \bigwedge_{i \in I} x_{i, f(i)}$. Suppose $\ell \not\le r$;
by totality, $r < \ell$.
\emph{Case 1:} There exists $m \in C$ with $r < m < \ell$. For every $i \in I$, since
$m < \ell \le \bigvee_j x_{i,j}$, there exists $f(i) \in J(i)$ with $x_{i, f(i)} \not\le m$,
hence $m \le x_{i, f(i)}$ by totality. Thus $m \le \bigwedge_i x_{i, f(i)} \le r$,
contradicting $r < m$.
\emph{Case 2:} No $m \in C$ satisfies $r < m < \ell$. For every $i \in I$, since
$r < \ell \le \bigvee_j x_{i,j}$, there exists $f(i) \in J(i)$ with $x_{i, f(i)} \not\le r$.
If $\ell \not\le x_{i, f(i)}$, totality would yield $r < x_{i, f(i)} < \ell$, contradicting
the case hypothesis; hence $\ell \le x_{i, f(i)}$ for all $i \in I$, giving
$\ell \le \bigwedge_i x_{i, f(i)} \le r$, contradiction. Thus $\ell \le r$, and
Proposition~\ref{prop:self-duality} completes~(1).

(2) In $\prod_{\alpha \in A} L_\alpha$, arbitrary infima and suprema are computed
coordinatewise, so $(\CDinf)$ holds at each coordinate $\alpha \in A$ and hence in the
product.

(3) If $\varphi : L \hookrightarrow P$ is an injective map preserving arbitrary infima
and suprema into a completely distributive complete lattice $P = \prod_{\alpha \in A} C_\alpha$,
then $\varphi(\bigwedge_i \bigvee_j x_{i,j}) = \bigwedge_i \bigvee_j \varphi(x_{i,j})
= \bigvee_f \bigwedge_i \varphi(x_{i, f(i)}) = \varphi(\bigvee_f \bigwedge_i x_{i, f(i)})$
by Proposition~\ref{prop:subset-collapse}, and injectivity of $\varphi$ yields equality in~$L$.
\end{proof}

Raney~\cite{Raney1953} proved the converse of Theorem~\ref{thm:raney-pos}(3): every
completely distributive complete lattice embeds as a complete sublattice into a direct
product of complete chains. By contrast, the \emph{literal} wording of Clause~1---``generated
by complete chain products''---fails under both standard algebraic notions of generation:

\begin{theorem}[Failure of literal Clause 1 on $M_3$]
\label{thm:clause1-literal}
Let $M_3 = \{\bot, a, b, c, \top\}$ be the $5$-element diamond lattice (in which $a, b, c$
form a $3$-element antichain with $\bot < a, b, c < \top$), and let $C_2 = \{0 < 1\}$ and
$C_3 = \{0 < 1 < 2\}$ denote the $2$- and $3$-element complete chains.
\begin{enumerate}
  \item \emph{(Internal generation by complete chain products)} For each $x \in \{a, b, c\}$,
    the subset $C^{(x)} := \{\bot, x, \top\} \subseteq M_3$ is a complete sublattice of
    $M_3$ isomorphic to the $1$-factor complete chain product $C_3^{\mathbf{1}}$, and
    $M_3 = C^{(a)} \cup C^{(b)} \cup C^{(c)}$ is internally generated by them, yet $M_3$
    is not completely distributive.
  \item \emph{(Complete-join quotient generation from a complete chain product)} Define
    $\pi : C_2^3 \to M_3$ by $\pi(0,0,0) = \bot$, $\pi(1,0,0) = a$, $\pi(0,1,0) = b$,
    $\pi(0,0,1) = c$, and $\pi(v) = \top$ whenever $v \in C_2^3$ has at least two coordinates
    equal to~$1$. Then $\pi$ is surjective and preserves arbitrary suprema
    $\pi(\bigvee S) = \bigvee \pi[S]$ for all $S \subseteq C_2^3$, yet $M_3$ is not
    completely distributive.
\end{enumerate}
\end{theorem}

\begin{proof}
In $M_3$, $a \wedge (b \vee c) = a \wedge \top = a \neq \bot = \bot \vee \bot = (a \wedge b) \vee (a \wedge c)$,
so $M_3$ fails binary distributivity and hence fails complete distributivity by
Proposition~\ref{prop:subset-collapse}. For~(1), any subset $S \subseteq C^{(x)}$ has
$\bigwedge_{M_3} S \in \{\bot, x, \top\} = C^{(x)}$ and $\bigvee_{M_3} S \in C^{(x)}$
(including $\bigwedge \emptyset = \top$ and $\bigvee \emptyset = \bot$), so each $C^{(x)}$
embeds $C_3^{\mathbf{1}}$ as a complete sublattice of $M_3$, and their union is all of $M_3$.
For~(2), define $\rho : M_3 \to C_2^3$ by $\rho(\bot) = (0,0,0)$, $\rho(a) = (1,0,0)$,
$\rho(b) = (0,1,0)$, $\rho(c) = (0,0,1)$, and $\rho(\top) = (1,1,1)$. Then $\pi(\rho(y)) = y$
for all $y \in M_3$ (proving surjectivity), and direct inspection of the five values of $y$
shows the Galois connection $\pi(v) \le y \iff v \le \rho(y)$ for all $v \in C_2^3$ and
$y \in M_3$. As a lower adjoint in a Galois connection between complete lattices, $\pi$
preserves all suprema: for every $S \subseteq C_2^3$ and $z \in M_3$,
$\pi(\bigvee S) \le z \iff \bigvee S \le \rho(z) \iff (\forall v \in S,\, v \le \rho(z))
\iff (\forall v \in S,\, \pi(v) \le z) \iff \bigvee \pi[S] \le z$.
\end{proof}

\section{Sharpness of Cardinality 5: All Lattices with \texorpdfstring{$|L| \le 4$}{|L| <= 4} Are Completely Distributive}

\begin{lemma}[Five-Distinct-Elements Lemma]
\label{lem:five-distinct}
Let $L$ be any complete lattice. If $x, y, z \in L$ satisfy
$x \wedge (y \vee z) \neq (x \wedge y) \vee (x \wedge z)$, then the five elements
\[
  y \wedge z, \quad y, \quad z, \quad y \vee z, \quad x
\]
are pairwise distinct in $L$.
\end{lemma}

\begin{proof}
First, $y \not\le z$: for if $y \le z$, then $y \vee z = z$ and $x \wedge y \le x \wedge z$,
giving $(x \wedge y) \vee (x \wedge z) = x \wedge z = x \wedge (y \vee z)$, a contradiction.
Symmetrically, $z \not\le y$. From $y \not\le z$ and $z \not\le y$, all six pairwise
inequalities among $\{y \wedge z, y, z, y \vee z\}$ follow immediately:
$y \wedge z = y \Rightarrow y \le z$; $y \wedge z = z \Rightarrow z \le y$;
$y \wedge z = y \vee z \Rightarrow y \le y \vee z = y \wedge z \le z$; $y = z \Rightarrow y \le z$;
$y = y \vee z \Rightarrow z \le y$; and $z = y \vee z \Rightarrow y \le z$.

Finally, $x$ cannot equal any of the four elements $w \in \{y \wedge z, y, z, y \vee z\}$,
because for each such $w$ the distributive identity $w \wedge (y \vee z) = (w \wedge y) \vee (w \wedge z)$
holds identically by absorption and idempotence:
\begin{align*}
  (y \wedge z) \wedge (y \vee z) &= y \wedge z = (y \wedge z) \vee (y \wedge z)
    = ((y \wedge z) \wedge y) \vee ((y \wedge z) \wedge z), \\
  y \wedge (y \vee z) &= y = y \vee (y \wedge z) = (y \wedge y) \vee (y \wedge z), \\
  z \wedge (y \vee z) &= z = (z \wedge y) \vee z = (z \wedge y) \vee (z \wedge z), \\
  (y \vee z) \wedge (y \vee z) &= y \vee z = ((y \vee z) \wedge y) \vee ((y \vee z) \wedge z).
\end{align*}
Hence $(y \wedge z, y, z, y \vee z, x)$ defines an injective map $\{0,1,2,3,4\} \hookrightarrow L$.
\end{proof}

\begin{theorem}[Complete distributivity for $|L| \le 4$]
\label{thm:size4}
Every complete lattice $L$ of cardinality $|L| \le 4$ is completely distributive.
\end{theorem}

\begin{proof}
Since $|L| \le 4$, $L$ is finite. Suppose $L$ were not completely distributive; by
Theorem~\ref{thm:finite-cd}, $L$ would fail binary distributivity, so Lemma~\ref{lem:five-distinct}
would produce an injection $\{0,1,2,3,4\} \hookrightarrow L$, contradicting $|L| \le 4$ by the
pigeonhole principle. (Concretely, up to isomorphism there are $1 + 1 + 1 + 2 = 5$ complete
lattices of size $1 \le |L| \le 4$: the chains $C_1, C_2, C_3, C_4$ and the $4$-element
Boolean diamond $B_2 = C_2 \times C_2$, all of which are completely distributive.)
\end{proof}

\section{Unconditional Refutation of Clause 2 and Master Disproof}

At the minimal cardinality $|L| = 5$ permitted by Theorem~\ref{thm:size4}, there exist
\emph{two} non-isomorphic complete lattices failing complete distributivity: the diamond
lattice $M_3 = \{\bot, a, b, c, \top\}$ and the pentagon lattice
$N_5 = \{\bot, u, v, w, \top\}$, whose partial order is defined by $\bot < u < v < \top$
and $\bot < w < \top$, with $w$ incomparable to both $u$ and $v$ (Figure~\ref{fig:hasse}).

\begin{figure}[htbp]
\centering
\setlength{\unitlength}{1pt}
\begin{tabular}{ccc}
\begin{picture}(120,105)(-10,-10)
  \put(50,0){\circle*{5}}
  \put(10,42){\circle*{5}}
  \put(50,42){\circle*{5}}
  \put(90,42){\circle*{5}}
  \put(50,84){\circle*{5}}
  \put(50,0){\line(-1,1){40}}
  \put(50,0){\line(0,1){84}}
  \put(50,0){\line(1,1){40}}
  \put(50,84){\line(-1,-1){40}}
  \put(50,84){\line(1,-1){40}}
  \put(46,-11){\small $\bot$}
  \put(-1,39){\small $a$}
  \put(55,39){\small $b$}
  \put(96,39){\small $c$}
  \put(46,89){\small $\top$}
\end{picture}
& \hspace{28pt} &
\begin{picture}(120,105)(-10,-10)
  \put(50,0){\circle*{5}}
  \put(22,28){\circle*{5}}
  \put(22,56){\circle*{5}}
  \put(78,42){\circle*{5}}
  \put(50,84){\circle*{5}}
  \put(50,0){\line(-1,1){28}}
  \put(22,28){\line(0,1){28}}
  \put(22,56){\line(1,1){28}}
  \put(50,0){\line(2,3){28}}
  \put(78,42){\line(-2,3){28}}
  \put(46,-11){\small $\bot$}
  \put(9,25){\small $u$}
  \put(9,53){\small $v$}
  \put(84,39){\small $w$}
  \put(46,89){\small $\top$}
\end{picture}
\\[4pt]
$M_3 \cong M_3^{\op}$ ($|M_3|=5$, modular) && $N_5 \cong N_5^{\op}$ ($|N_5|=5$, non-modular)
\end{tabular}
\caption{Hasse diagrams of the two non-isomorphic $5$-element complete lattices failing
complete distributivity.}
\label{fig:hasse}
\end{figure}

Recall that a lattice $L$ is \emph{modular} if $x \le z \implies x \vee (y \wedge z) = (x \vee y) \wedge z$
for all $x, y, z \in L$. Modularity is preserved under order isomorphism.

\begin{proposition}[Properties and non-isomorphism of $M_3$, $M_3^{\op}$, and $N_5$]
\label{prop:M3-N5}
\begin{enumerate}
  \item Both $M_3$ and $N_5$ are $5$-element complete lattices, and $M_3$ is self-dual
    ($M_3^{\op} \cong M_3$, swapping $\bot$ and $\top$ while fixing $a, b, c$).
  \item Both $M_3$ (hence $M_3^{\op}$) and $N_5$ fail complete distributivity: in $N_5$,
    \[
      v \wedge (u \vee w) = v \wedge \top = v \;\neq\; u = u \vee \bot = (v \wedge u) \vee (v \wedge w).
    \]
  \item $M_3$ is modular, whereas $N_5$ is not modular: $u \le v$, yet
    $u \vee (w \wedge v) = u \vee \bot = u \neq v = \top \wedge v = (u \vee w) \wedge v$.
    Consequently, $M_3 \not\cong N_5$ and $M_3^{\op} \not\cong N_5$.
    (Independently, $M_3$ has $3$ atoms and maximal chain size $3$, whereas $N_5$ has $2$
    atoms and maximal chain size $4$.)
\end{enumerate}
\end{proposition}

\begin{proof}
Items~(1) and~(2) follow directly from the defining order relations and
Proposition~\ref{prop:subset-collapse}. For~(3), if $x \le z$ in $M_3$, then either
$x = \bot$ (giving $\bot \vee (y \wedge z) = y \wedge z = (\bot \vee y) \wedge z$),
$z = \top$ (giving $x \vee (y \wedge \top) = x \vee y = (x \vee y) \wedge \top$), or
$x = z \in \{a, b, c\}$ (giving $x \vee (y \wedge x) = x = (x \vee y) \wedge x$). Thus
$M_3$ is modular. Since $N_5$ fails modularity on $(u, w, v)$ and order isomorphisms
preserve $\le$, $\wedge$, $\vee$, and hence modularity, no order isomorphism exists
between $M_3$ (or $M_3^{\op} \cong M_3$) and $N_5$.
\end{proof}

We also check the alternative standard order-theoretic reading of ``the dual of the
three-element antichain construction'' via the down-set (order-ideal) completion:

\begin{proposition}[Dual of the down-set completion of $A_3$]
\label{prop:downset-A3}
Let $A_3$ be the $3$-element antichain. Every subset of $A_3$ is a down-set, so its
down-set completion $\Down(A_3) = (\mathcal{P}(A_3), \subseteq) \cong C_2^3$ and its
order dual $\Down(A_3)^{\op} \cong (C_2^3)^{\op} = B_3^{\op}$ have cardinality $8$ and
\emph{satisfy} complete distributivity.
\end{proposition}

\begin{proof}
$C_2$ is a $2$-element complete chain, so $C_2^3$ is completely distributive by
Theorem~\ref{thm:raney-pos}(2), and its dual $(C_2^3)^{\op}$ is completely distributive by
Proposition~\ref{prop:self-duality}, with $|(C_2^3)^{\op}| = 2^3 = 8$.
\end{proof}

Finally, we prove a construction-independent impossibility theorem refuting Clause~2
for \emph{every} candidate complete lattice $L_*$, under both cardinality-minimality (with
uniqueness up to order isomorphism) and order/sublattice-embedding minimality:

\begin{theorem}[Universal refutation of Clause 2]
\label{thm:clause2-false}
Let $L_*$ be any complete lattice.
\begin{enumerate}
  \item \emph{(Cardinality minimality + uniqueness)} It is \textbf{false} that $L_*$ fails
    complete distributivity, injects into every complete lattice $M$ failing complete
    distributivity ($|L_*| \le |M|$), and is unique up to order isomorphism among complete
    lattices of that minimum cardinality failing complete distributivity.
  \item \emph{(Order / sublattice embedding minimality)} It is \textbf{false} that $L_*$
    fails complete distributivity and order-embeds ($x \le y \iff e(x) \le e(y)$, which
    subsumes sublattice and complete-sublattice embeddings) into every complete lattice
    $M$ failing complete distributivity.
\end{enumerate}
\end{theorem}

\begin{proof}
(1) Suppose $L_*$ satisfies those hypotheses. Applying $|L_*| \le |M|$ to $M = M_3$ yields
an injection $L_* \hookrightarrow M_3 \cong \{0,1,2,3,4\}$. By Theorem~\ref{thm:size4},
$L_*$ cannot inject into $\{0,1,2,3\}$ since $L_*$ fails complete distributivity; hence
the injection $L_* \hookrightarrow \{0,1,2,3,4\}$ must be surjective (missing a point $k$
would compress into $\{0,1,2,3\}$), so $|L_*| = 5 = |M_3| = |N_5|$. By uniqueness at
minimum cardinality applied to $M_3$ and $N_5$, we obtain order isomorphisms
$L_* \cong M_3$ and $L_* \cong N_5$, whence $M_3 \cong N_5$, contradicting
Proposition~\ref{prop:M3-N5}(3).

(2) Suppose $L_*$ fails complete distributivity and admits order-embeddings
$e_1 : L_* \hookrightarrow_o M_3$ and $e_2 : L_* \hookrightarrow_o N_5$. Any
order-embedding is injective by antisymmetry. Since $L_*$ cannot inject into $\{0,1,2,3\}$
by Theorem~\ref{thm:size4}, both $e_1$ and $e_2$ are bijective order-embeddings onto
$5$-element lattices, hence order isomorphisms $L_* \cong M_3$ and $L_* \cong N_5$,
again yielding $M_3 \cong N_5$, contradiction.
\end{proof}

\begin{theorem}[Master Disproof of Conjecture \texttt{00000008559}]
\label{thm:master-disproof}
For every proposition $P_1$ representing Clause~1 (whether the literal internal generation
equivalence, the complete-join quotient equivalence, or Raney's complete-sublattice
equivalence) and every complete lattice $L_*$ representing ``the dual of the three-element
antichain construction'' (including $M_3^{\op}$ and $\Down(A_3)^{\op}$), the conjunction
$P_1 \wedge P_2(L_*)$ of Conjecture \texttt{00000008559} is \textbf{false}.
\end{theorem}

\begin{proof}
Immediate from Theorem~\ref{thm:clause2-false}.
\end{proof}

\section{Formal Verification in Core Lean 4 and Computational Audit}

All definitions and theorems in Sections~2--5 are formalized from scratch in Core Lean 4
(\texttt{leanprover/lean4:v4.19.0}) in \texttt{lean/Main.lean}, compiled with
\texttt{-DwarningAsError=true} and zero external dependencies. Every major declaration is
checked with \texttt{\#print axioms}, confirming reliance solely on Lean's standard
foundational axioms (\texttt{propext}, \texttt{Classical.choice}, \texttt{Quot.sound}),
with zero \texttt{sorryAx} and zero \texttt{Lean.ofReduceBool} (\texttt{native\_decide}).
Table~\ref{tab:lean-map} records the correspondence between this paper and \texttt{lean/Main.lean}.

\begin{table}[htbp]
\centering
\footnotesize
\renewcommand{\arraystretch}{1.2}
\begin{tabular}{@{}>{\raggedright\arraybackslash}p{0.41\textwidth}>{\raggedright\arraybackslash}p{0.55\textwidth}@{}}
\toprule
\textbf{Mathematical Result} & \textbf{Lean 4 Declaration (\texttt{Conjecture00000008559.*})} \\
\midrule
Definition~\ref{def:complete-lattice} (Complete lattice) & \texttt{CompleteLattice} \\
Definition~\ref{def:cd} ($\CDinf$, $\CDsup$, CD) & \texttt{CDInfSup}, \texttt{CDSupInf}, \texttt{IsCompletelyDistributive} \\
Proposition~\ref{prop:self-duality} (Raney self-duality) & \texttt{cdInfSup\_iff\_cdSupInf}, \texttt{isCompletelyDistributive\_dual\_iff} \\
Proposition~\ref{prop:subset-collapse} (Universe invariance) & \texttt{IsCompletelyDistributive.iInf\_iSup\_univ} \\
Theorem~\ref{thm:finite-cd} (Finite CD $\iff$ distributive) & \texttt{finite\_completelyDistributive\_iff\_distributive} \\
Theorem~\ref{thm:raney-pos} (Chain-product sufficiency) & \texttt{isCompletelyDistributive\_of\_embedsInChainProduct} \\
Theorem~\ref{thm:clause1-literal}(1) (Literal internal false) & \texttt{clause1\_literal\_internal\_false} \\
Theorem~\ref{thm:clause1-literal}(2) (Literal quotient false) & \texttt{clause1\_literal\_joinQuotient\_false} \\
Lemma~\ref{lem:five-distinct} ($5$ distinct elements) & \texttt{five\_distinct\_of\_not\_distributive} \\
Theorem~\ref{thm:size4} ($|L| \le 4 \implies$ CD) & \texttt{completelyDistributive\_of\_injective\_fin\_le\_four} \\
Proposition~\ref{prop:M3-N5} ($M_3$, $N_5$, $M_3^{\op} \not\cong N_5$) & \texttt{not\_isCompletelyDistributive\_M3}, \texttt{not\_isCompletelyDistributive\_N5}, \texttt{orderIsoDualM3}, \texttt{not\_orderIso\_M3\_N5}, \texttt{not\_orderIso\_dualM3\_N5} \\
Proposition~\ref{prop:downset-A3} ($\Down(A_3)^{\op}$ size $8$, CD) & \texttt{equivDualDownSetA3Fin8}, \texttt{isCompletelyDistributive\_DualDownSetA3} \\
Theorem~\ref{thm:clause2-false} (Universal Clause 2 disproof) & \texttt{clause2\_false\_by\_cardinality}, \texttt{clause2\_false\_by\_embedding} \\
Theorem~\ref{thm:master-disproof} (Master Disproof) & \texttt{conjecture\_00000008559\_disproved} \\
\bottomrule
\end{tabular}
\caption{Exact correspondence between mathematical statements and Core Lean 4 declarations.}
\label{tab:lean-map}
\end{table}

In addition, the standalone Python script \texttt{verify.py} exhaustively enumerates all
unlabeled complete lattices of cardinalities $1 \le |L| \le 5$ (reproducing OEIS A006966:
$1, 1, 1, 2, 5$), confirms $0$ non-completely-distributive complete lattices for $|L| \le 4$
and exactly $2$ non-isomorphic counterexamples ($M_3$ and $N_5$) at $|L| = 5$, verifies
$|\Down(A_3)^{\op}| = 8$ with complete distributivity holding, and checks $\pi(\bigvee S) = \bigvee \pi[S]$
across all $2^8 = 256$ subsets $S \subseteq C_2^3$ for the complete-join quotient $\pi : C_2^3 \twoheadrightarrow M_3$.

\begin{thebibliography}{9}

\bibitem{Birkhoff1967}
G.~Birkhoff,
\emph{Lattice Theory}, 3rd ed.,
American Mathematical Society Colloquium Publications, Vol.~25, Providence, RI, 1967.

\bibitem{DaveyPriestley2002}
B.~A.~Davey and H.~A.~Priestley,
\emph{Introduction to Lattices and Order}, 2nd ed.,
Cambridge University Press, Cambridge, 2002.

\bibitem{Gierz2003}
G.~Gierz, K.~H.~Hofmann, K.~Keimel, J.~D.~Lawson, M.~Mislove, and D.~S.~Scott,
\emph{Continuous Lattices and Domains},
Encyclopedia of Mathematics and its Applications, Vol.~93, Cambridge University Press, Cambridge, 2003.

\bibitem{Raney1952}
G.~N.~Raney,
Completely distributive complete lattices,
\emph{Proceedings of the American Mathematical Society} \textbf{3} (1952), 677--680.

\bibitem{Raney1953}
G.~N.~Raney,
A subdirect-union representation for completely distributive complete lattices,
\emph{Proceedings of the American Mathematical Society} \textbf{4} (1953), 518--522.

\end{thebibliography}

\end{document}
